\input{noncyclic/stabilizer_multiplicity_criterion_root.tex}

\begin{proposition}\label{prop:octic-no318}
No elliptic curve over a number field of degree $8$ has a subgroup
isomorphic to $C_3\times C_{18}$.
\end{proposition}
\begin{proof}
Put $F=\mathbb Q(\zeta_3)$ and $X=X_1(3,18)_F$.
The Weil pairing forces $F$ into a proposed field of definition,
so the subgroup gives an effective noncuspidal $F$-rational
trace divisor $D$ of degree $4$ on $X$.
We have $g(X)=10$ and $\operatorname{rank}J(F)=0$ by
\cite[Theorem~4.1]{DS}. The gonality bound in \cite{NV}
gives $\operatorname{gon}_{\mathbb F_7}X\geq6$,
and hence $\operatorname{gon}_F X\geq6$.

We use the split primes $7$ and $13$.
Both generate $\operatorname{Gal}(\mathbb Q(\zeta_{18})/F)$.
At either prime the cusp signature is $1^{18},3^6$, and
each closed residue cusp therefore lifts to an $F$-rational
cuspidal divisor. Entirely cuspidal reduction of $D$ at either
prime contradicts the odd-prime cuspidal specialization criterion.
The noncuspidal closed-place counts in degrees one through four are
\[
 \begin{array}{c|rrrr}
 p&1&2&3&4\\\hline
 7&0&27&144&513\\
 13&0&54&846&7506.
 \end{array}
\]
Newton identities applied to the Frobenius polynomial of
$J_H(54)$, $H=\langle19\rangle$, independently check these counts.

There are $27$ $F$-automorphisms
\[
 (P,Q)\longmapsto(aP+6bQ,cP+dQ),\quad
 b,c\in\mathbb Z/3\mathbb Z,\quad d\in\{1,5,7\},
 \quad a=d^{-1}\pmod3.
\]
They preserve the chosen Weil pairing and commute with the Hecke
operators at primes congruent to $1$ modulo $3$.
Their action on geometric noncuspidal points is free, as proved
after Lemma~\ref{lem:octic-stabilizer-multiplicity}.

For a split prime $q\nmid18p$ write
$A_q=T_q-1-q\langle q\rangle$.
Eichler--Shimura and odd unramified torsion specialization show
that $A_q$ annihilates $J(F)$; its divisor formula vanishes on
every $F$-rational cusp. At $7$ the $27$, $144$, and $513$
noncuspidal places give respectively $1$, $8$, and $19$
automorphism orbits.
The exact $A_{13}$ and $A_{19}$ tests exclude the quadratic orbit,
four cubic orbits, and eleven quartic orbits.
The four remaining cubic orbits have size $9$ and are excluded
by Lemma~\ref{lem:octic-stabilizer-multiplicity}, since
$4+1<6$ and $3\nmid4$.
For two quadratic supports we fix the first by transitivity and
test all $27$ possibilities for the second, including repetition;
all are excluded. Thus $\overline D_7$ must be one of the $216$
quartic places in eight remaining orbits.

We next compare the two primes without assuming that cusps generate
$J(F)$. Choose a rational cusp $c_0$ and let $C_F\subset J(F)$
be generated by the $24$ global cusp-block classes
$[B-\deg(B)c_0]$.
The blocks are identified at both primes by valuation vectors
of $65$ integral modular units: $a_3$, $a_1^3-27a_3$,
the coordinates of $kQ+jP$ for $1\leq k\leq9$,
$0\leq j\leq2$, and $x(kQ)-x(kQ+P)$.
These units have no noncuspidal zeros or poles and distinguish
all $24$ blocks. The proved finite class-group calculations give
\[
 \begin{split}
 J(\mathbb F_7)&\simeq C_3^3\times C_9^2
                  \times C_{63}^2\times C_{189},\\
 C_F&\simeq C_3^4\times C_{63}^3,\\
 J(\mathbb F_7)/\overline C_F&\simeq C_3\times C_3\times C_9.
 \end{split}
\]
Here $|C_F|=20253807$. Injectivity of torsion specialization
proves $9J(F)\subset C_F$.

Let $t=[D-4c_0]$.
The pair $(9\overline t_7,9\overline t_{13})$ must lie in the
diagonal image of $C_F$. We compute this image using the actual
relation matrices and the paired cusp-class vectors.
The $216$ source quartics have $135$ distinct nine-multiple classes.
At $13$ we test all $54\binom{19}{2}=9234$ quadratic-plus-cusp
cases, $846\cdot18=15228$ cubic-plus-cusp cases,
$\binom{55}{2}=1485$ quadratic pairs, and $278$ quartic
automorphism representatives: $26225$ cases in total.
Reducing quartics by automorphisms is valid because the complete
source set and the diagonal cusp subgroup are stable under the
global action. On the nine-multiple class the change of base cusp
is the global cuspidal translation $36[g(c_0)-c_0]$.

Exact integer Hermite normal form leaves $270$ cubic-plus-cusp
cases supported on $54$ distinct cubics, and nine quartic
representatives. The $54$ cubics form six orbits of size $9$;
Lemma~\ref{lem:octic-stabilizer-multiplicity} excludes all of them.
For each of the nine quartic representatives at $13$, the divisor
$A_7x$ is nonprincipal. There are no remaining reductions.

The certificate independently replays the finite-field arithmetic,
all automorphism partitions, the cusp-block comparison, the
integer-lattice intersection, and the assembly of every Hecke
decision. Native isogeny enumerations agree exactly with the
fast $T_{13}$ divisor of degree $28$ at $7$ and the fast $T_7$
divisor of degree $32$ at $13$.
The class groups are proved before installing the direct forward
map. Every class computed by this map is checked through
$E=D_0+R+\operatorname{div}(b)$ in both maximal orders, with
$E$ and $R$ supported on the native factor base.
A fresh complete calculation at $7$ repeats the two-prime
comparison in an independently computed class-group presentation.
Hence the proposed subgroup does not exist.
\end{proof}
